%
% File: abstract.tex
% Author: Solal Rapaport
% Description: Thesis abstract — English then French.

\chapter*{Abstract}
\begin{SingleSpace}
\initial{P}ublic software repositories are the foundation on which the modern software supply chain is built.
Package managers, continuous integration pipelines, and reproducible-build systems all resolve dependencies against repository state under an implicit assumption: that a named version label like \texttt{v1.0} designates stable, fixed content.
Git's design, however, separates immutable content-addressed objects from mutable references.
Tags, branches, and other symbolic names can be silently updated at any time, while the version string in a downstream configuration file remains unchanged.
This gap between assumption and reality is the subject of this thesis.

The thesis measures that gap empirically, at ecosystem scale, using \textbf{Software Heritage} --- the largest public archive of source code --- as its primary data source.
By comparing consecutive archival snapshots of the same repository, it recovers the before-and-after states of alterations that would otherwise be irrecoverable from the forge once the repository has been updated.
This snapshot-pair methodology is applied to hundreds of millions of repositories across three empirical studies.

The \textbf{first study} analyses tag alterations in 45.5 million repositories observed over nine years.
It detects \num{10.2} million alterations in \num{189000} repositories (\num{0.41}\% of all repositories with tags), produces a taxonomy of alteration forms distinguishing tag moves that changed content, metadata-only moves, and deletions, and demonstrates through a cross-analysis with the Nixpkgs package set that tag moves cause concrete build reproducibility failures in real downstream package ecosystems.

The \textbf{second study} examines commit-graph history alterations in \num{111} million repositories, identifying \num{8.7} million root-cause rewrites across \num{1.22} million repositories (\num{1.10}\%).
It characterises the distribution of alterations by branch type and repository popularity, introduces a taxonomy of alteration forms (directory-content changes, metadata-only changes, and branch relocations), and documents two security-relevant patterns at ecosystem scale: \num{13} million removals of files whose names suggested secrets across \num{75000} repositories, and retroactive license changes in \num{32000} repositories.
The study further produces \textit{GitHistorian}, a command-line tool that operationalises history alteration detection on the Software Heritage graph for practitioner use.

The \textbf{third study} pursues the question left open by the second: of the removed files that had contained credentials and that remain archived in Software Heritage, how many contain detectable, syntactically valid secrets?
Working from the corpus of removals of files whose names suggested secrets, it applies a multi-stage pipeline that recovers archived blobs, scans them with two independent secret detectors, and links findings back to their repository provenance.
Among \num{281728} downloaded blobs, \num{9.18}\% yield at least one finding.
Among private-key blobs specifically, \num{86.4}\% of flagged blobs pass offline syntactic validation.
Whether these credentials remain active is the central open question: the study suggests a detection infrastructure and a staged responsible-disclosure protocol designed to answer it.

Together, the three studies establish that repository mutability is an active, recurring phenomenon with measurable supply-chain consequences, and that without archival snapshots all we can observe in a public repository is what its current references let us see.
\end{SingleSpace}
\clearpage

\chapter*{Résumé}
\begin{SingleSpace}
\initial{L}es dépôts logiciels publics constituent la fondation sur laquelle repose la chaîne d'approvisionnement logicielle moderne.
Les gestionnaires de paquets, les pipelines d'intégration continue et les systèmes de construction reproductibles résolvent tous leurs dépendances en s'appuyant sur l'état des dépôts, sous l'hypothèse implicite qu'une étiquette de version comme \texttt{v1.0} désigne un contenu stable et fixe.
La conception de Git sépare cependant les objets immuables adressés par contenu des références mutables~: les tags, les branches et autres noms symboliques peuvent être mis à jour silencieusement à tout moment, tandis que la chaîne de version dans les fichiers de configuration des consommateurs reste inchangée.
Ce fossé entre l'hypothèse et la réalité est le sujet de cette thèse.

La thèse mesure cet écart empiriquement, à l'échelle de l'écosystème, en utilisant \textbf{Software Heritage} --- la plus grande archive publique de code source --- comme principale source de données.
En comparant des snapshots archivés successifs d'un même dépôt, elle rend observables les états avant et après les altérations qui seraient autrement irrecouvrables depuis la forge une fois le dépôt mis à jour.
Cette méthodologie par paires d'snapshots est appliquée à des centaines de millions de dépôts, au travers de trois études empiriques.

La \textbf{première étude} analyse les altérations de tags dans \num{45,5}~millions de dépôts observés sur neuf ans.
Elle détecte \num{10,2}~millions d'altérations dans \num{189000}~dépôts (\num{0,41}\,\% de tous les dépôts avec au moins un tag), produit une taxonomie distinguant les déplacements de tags modifiant le contenu, les déplacements ne modifiant que les métadonnées et les suppressions, et démontre, par une analyse croisée avec l'ensemble de paquets Nixpkgs, que les déplacements de tags provoquent des échecs concrets de reproductibilité dans des écosystèmes réels.

La \textbf{deuxième étude} examine les altérations de l'historique de commits dans \num{111}~millions de dépôts, identifiant \num{8,7}~millions de commits directement ciblés par une alteration dans \num{1,22}~million de dépôts (\num{1,10}\,\%).
Elle caractérise la distribution des altérations par type de branche et par popularité du dépôt, introduit une taxonomie des formes d'altération, et documente deux schémas à caractère sécuritaire à l'échelle de l'écosystème~: \num{13}~millions de suppressions de fichiers à caractère secret dans \num{75000}~dépôts, et des modifications rétroactives de licences dans \num{32000}~dépôts.
L'étude produit également \textit{GitHistorian}, un outil en ligne de commande qui opérationnalise la détection d'altérations d'historique sur le graphe Software Heritage.

La \textbf{troisième étude} poursuit la question laissée ouverte par la seconde~: parmi les fichiers supprimés contenant des informations secrètes qui restent archivés dans Software Heritage, combien contiennent des secrets détectables et syntaxiquement valides~?
Parmi \num{281728}~blobs téléchargés, \num{9,18}\,\% sont détectés comme ayant au moins un secret~; parmi ces blobs, \num{86,4}\,\% passent la validation syntaxique hors ligne.
La thèse propose une infrastructure de détection et un protocole de divulgation responsable qui permettrait de tester la validité des secrets pour répondre à cette question.

Ensemble, les trois études établissent que la mutabilité des dépôts est un phénomène actif et récurrent aux conséquences mesurables sur la chaîne d'approvisionnement, et que, sans snapshots archivés, tout ce que nous pouvons observer dans un dépôt public est ce que ses références actuelles nous laissent voir.
\end{SingleSpace}
\clearpage
